\documentclass[10pt,a4paper]{amsart}
\usepackage[margin=19mm]{geometry}
\usepackage{amsmath,amssymb,mathtools}
\usepackage[scr=rsfs,cal=euler]{mathalfa}
\usepackage{xfrac}
\usepackage[unicode,hidelinks,bookmarks=false]{hyperref}
\newcommand{\ie}{i.e.\ }
\newcommand{\cf}{cf.\ }
\newcommand{\resp}{respectively\ }
\newcommand{\Subsection}{\subsection}
\providecommand{\nobreakdash}{\nobreak\hbox{-}}
\setlength{\emergencystretch}{2em}
\multlinegap0pt
\usepackage{tikz}
\newtheorem{theorem}{Theorem}
\title{Group Theory of the Kolakoski Sequence}
\author{Noah MacAulay}
\date{Source: arXiv:2605.14234v1 (CC BY 4.0)}
\begin{document}
\maketitle

\noindent\emph{Source and licence.} Group Theory of the Kolakoski Sequence. Version arXiv:2605.14234v1 (CC BY 4.0), distributed by arXiv under the Creative Commons Attribution 4.0 International licence, http://creativecommons.org/licenses/by/4.0/, as recorded in the arXiv licence metadata for this exact version and shown on the abstract page https://arxiv.org/abs/2605.14234v1. Full licence notice: \texttt{licenses/arxiv-2605.14234-CC-BY-4.0.txt}.\par\smallskip

\noindent\textit{Editorial scope.} Counted unit: the article's theorem theorem (source0.tex lines 477--488). The article's complete original proof of it is delivered with it (source0.tex lines 490--555, one \texttt{proof} environment of 66 lines). No mathematics is written, repaired or re-derived: the statement and the proof are the article's own lines, in the article's own notation, and no retained line is substituted or re-flowed.  Retained uncounted as the source-local context the proof needs: lines 469--471.  Nothing was omitted from any retained span, so the package carries no omission record.  No retained line contains a citation, so the wrapper carries no bibliography and no bibliography entry.  No retained line contains a figure environment, an image-inclusion command or drawing code, so no image is delivered and the package declares no asset dependency.  Editorial formatting only, in addition: an AMS \texttt{amsart} preamble that restates the article's own declarations and macros used by the retained lines, namely the environments \texttt{theorem} on separate counters, together with the standard packages those lines need.  The retained environments print in this document as Theorem 1 in order of appearance, whereas the article prints the same items with the numbering of its own full document.  The complete native source of the article as distributed in the arXiv e-print for this exact version is bundled unchanged as \texttt{sources/200568/source0.tex}.
\begin{equation}\label{defin}
    \psi_{n-2}(g_{LL}) - \psi_{n-2}(g_{LR}) = \psi_{n-2}(g_{RR}) - \psi_{n-2}(g_{RL}) + \delta(g_{Lf}=1)(p+q)(p-q)
\end{equation}

\begin{theorem}
    Using the definitions above, we have
    
    (a) $\mathcal{J}_n ^{p,q}$ is a group
    
    (b) $\psi_n$ is a homomorphism
    
    
    (c) $\Delta_n$ is a homomorphism
    
    
\end{theorem}

\begin{proof}
    We proceed by induction. The claims are clear for $n=1,2$. 
    
    \vspace{1em}
    
    (a) Proof $\mathcal{J}_n^{p,q}$ is a group: If we have two elements $g, h   \in \mathcal{J}_n^{p,q}$, then  $gh$ is either equal to $(g_f+ h_f, g_Lh_L,g_Rh_R)$ or $(g_f+ h_f, g_Lh_R,g_Rh_L)$, depending on whether $g_f$ is 0 or 1, respectively. In the first case we have $\Phi(\Delta(g_L h_L)) = \Phi(\Delta(g_L)) \Phi(\Delta(h_L)) =\Delta(g_R h_R)$, in the second $\Phi(\Delta(g_L h_R)) = \Phi(\Delta(g_L)) \Phi(\Delta(h_R)) = \Delta(g_R h_L)$ (because $\Phi$ is an involution). So $gh$ is in $\mathcal{J}_n^{p,q}$

    As for inverses, if $\Delta(g_L) = \Phi(\Delta(g_R))$ we have $\Delta(g_L^{-1}) = \Phi(\Delta(g_R^{-1}))$ and $\Delta(g_R^{-1}) = \Phi(\Delta(g_L^{-1}))$, this is enough to show $g^{-1} \in \mathcal{J}_n^{p,q}$

    \vspace{1em}
    
    (b.i) Proof $\psi_n$ is a homomorphism, $n$ odd: $$\psi_n (gh) = \frac{\psi_{n-1}((gh)_L) + \psi_{n-1}((gh)_R)}{p+q} = \frac{\psi_{n-1}(g_L) + \psi_{n-1}(g_R) + \psi_{n-1}(h_L) + \psi_{n-1}(h_R)}{p+q} = \psi_n(g) + \psi_n(h)$$

    (b.ii) Proof $\psi_n$ is a homomorphism, $n$ even. We consider two further sub-cases: $g_{Lf} = 0$ and $g_{Lf} = 1$
    
    (b.ii.i) $g_{Lf} = 0$. In this case we have

    $$\psi_n(gh) = \frac{2(\psi_{n-2}((gh)_{LL}) + \psi_{n-2}((gh)_{RL})) - \delta(h_{Lf})(p+q)(p-q) }{(p+q)^2}$$

    $$ = \frac{2(\psi_{n-2}(g_{LL}) + \psi_{n-2}(h_{LL}) +  \psi_{n-2}(g_{RL}) + \psi_{n-2}(h_{RL})) - \delta(h_{Lf})(p+q)(p-q) }{(p+q)^2}$$

    $$ = \psi_n(g) + \psi_n(h)$$

    (b.ii.ii) $g_{Lf} = 1$. In this case we have
    $$\psi_n(gh) = \frac{2(\psi_{n-2}((gh)_{LL}) + \psi_{n-2}((gh)_{RL})) - (1 - \delta(h_{Lf}))(p+q)(p-q) }{(p+q)^2}$$
    
    $$ = \frac{2(\psi_{n-2}(g_{LL}) + \psi_{n-2}(h_{LR}) +  \psi_{n-2}(g_{RL}) + \psi_{n-2}(h_{RR})) - (1 - \delta(h_{Lf}))(p+q)(p-q) }{(p+q)^2}$$
    By equation \ref{defin}, this is 
    $$  \frac{2(\psi_{n-2}(g_{LL}) + \psi_{n-2}(h_{LL}) +  \psi_{n-2}(g_{RL}) + \psi_{n-2}(h_{RL}) - \delta(h_{Lf})(p+q)(p-q)) - (1 - \delta(h_{Lf}))(p+q)(p-q)) }{(p+q)^2}$$

    $$ = \frac{2(\psi_{n-2}(g_{LL}) + \psi_{n-2}(h_{LL}) +  \psi_{n-2}(g_{RL}) + \psi_{n-2}(h_{RL})) - (1 + \delta(h_{Lf}))(p+q)(p-q)) }{(p+q)^2}$$

    $$= \psi_n(g) + \psi_n(h)$$
    
    \vspace{1em}

    (c.i) Proof $\Delta_n$ is a homomorphism. We consider subcases $n$ odd or even, and $g_f = 0$ or $1$. 

    (c.i.i) $n$ odd, $g_f = 0$. We have
    $$\Delta_n(gh) = ((gh)_f, \psi_{n-1}((gh)_L) - \psi_{n-1}((gh)_R))$$
    $$ = (h_f, \psi_{n-1}(g_L) - \psi_{n-1}(g_R) + \psi_{n-1}(h_L)  - \psi_{n-1}(h_R)) = \Delta_n(g)\Delta_n(h)$$

    (c.i.ii) $n$ odd, $g_f = 1$. We have 
    $$\Delta_n(gh) = ((gh)_f, \psi_{n-1}((gh)_L) - \psi_{n-1}((gh)_R))$$
    $$ = (1 + h_f, \psi_{n-1}(g_L) - \psi_{n-1}(g_R) - (\psi_{n-1}(h_L) - \psi_{n-1}(h_R)))$$
    $$ = \Delta_n(g)\Delta_n(h)$$

    (c.ii.ii). $n$ even, $g_f = 0$. We have 
    $$\Delta_n(gh) = ((gh)_f, (p+q)\psi_{n}(gh) - 2\psi_{n-1}((gh)_R))$$
    $$= (h_f, (p+q)(\psi_{n}(g) + \psi_n(h)) - 2(\psi_{n-1}(g_R) + \psi_{n-1}(h_R)))$$
    $$= \Delta_n(g)\Delta_n(h)$$

    (c.ii.ii) $n$ even, $g_f = 1$. We have
    $$\Delta_n(gh) = ((gh)_f, (p+q)\psi_{n}(gh) - 2\psi_{n-1}((gh)_R))$$
    \begin{equation}\label{lala} = (1+h_f, (p+q)(\psi_{n}(g) + \psi_n(h)) - 2(\psi_{n-1}(g_R) + \psi_{n-1}(h_L)))\end{equation}
    Now consider the sum $(p+q)\psi_n(h) - 2\psi_{n-1}(h_L)$. As $n$ is even this is equal to 
    $$\frac{2(\psi_{n-2}(h_{LL}) + \psi_{n-2}(h_{RL})) - \delta(h_{Lf})(p+q)(p-q)}{p+q} - 2\frac{\psi_{n-2}(h_{LL}) + \psi_{n-2}(h_{LR})}{p+q}$$
    $$=\frac{2(\psi_{n-2}(h_{RL}) - \psi_{n-2}(h_{LR})) - \delta(h_{Lf})(p+q)(p-q)}{p+q}$$
    By equation \ref{defin}, this is equal to
    $$\frac{2(\psi_{n-2}(h_{RR}) - \psi_{n-2}(h_{LL})) + \delta(h_{Lf})(p+q)(p-q)}{p+q}$$
    $$ = 2\psi_{n-1}(h_R) - (p+q)\psi_n(h)$$
    Substituting this into \ref{lala} we have 
    $$\Delta_n(gh) = (1+h_f, (p+q)(\psi_n(g) - \psi_n(h)) - 2(\psi_{n-1}(g_R) - \psi_{n-1}(h_R)))$$
    $$ = \Delta_n(g)\Delta_n(h)$$

\end{proof}

\end{document}
